% The exact proof source passed strict native-context checks and kernel audits.
% Ordinary module verification and declaration checks remain pending.

\section{Uniform floors for band metrics and mean trees}

\begin{theorem}[Uniform floor for symmetric band metrics]
    \label{thm:replica_symmetric_band_floor}
    \lean{TensorPower.ReplicaTransport.exists_symBandMetric_floor}
    \leanok
    \uses{lem:replica_symmetric_metric}
    Fix a finite tensor product $\mathcal V=\bigotimes_{v\in V}\C^{n_v}$,
    with $n_v>0$, and let $t>0$. The replica metrics use one common label
    weight padded to the original full dimension $\dim\mathcal V$.
    There exists $C\ge0$, chosen before the copy number and partition,
    such that every $k\ge0$ and every partition $\pi=(P,Y,F)$ satisfy
    \begin{align}
        (k+2)^{-C}I\le\widetilde A_k(\pi),
    \end{align}
    where $\widetilde A_k(\pi)=\Pi_kA_k(\pi)+(I-\Pi_k)$ is the existing
    identity extension of the symmetric band metric and
    $A_k(\pi)=(W_{P,k}^{-1}W_{F,k}^{-1}W_{Y,k})^2$.
    Empty parts, $k=0$, and an empty set of sites are included.
\end{theorem}

\begin{proof}\leanok
    \uses{thm:replica_symmetric_floor, lem:replica_symmetric_metric}
    Choose the exponent from Theorem~\ref{thm:replica_symmetric_floor}
    for the fixed system and $t$. Its identity-extension inequality is
    \begin{align}
        (k+2)^{-C}I\le\Pi_kA_k(\pi)\Pi_k+(I-\Pi_k).
    \end{align}
    The band metric commutes with $\Pi_k$, and $\Pi_k^2=\Pi_k$.
    Hence $\Pi_kA_k(\pi)\Pi_k=\Pi_kA_k(\pi)$, proving the assertion.
\end{proof}

\begin{theorem}[Uniform floor for mean trees of symmetric band metrics]
    \label{thm:replica_symmetric_mean_floor}
    \lean{TensorPower.ReplicaTransport.exists_meanTree_symBandMetric_floor}
    \leanok
    \uses{lem:replica_symmetric_metric, def:mean_tree}
    Fix the same finite tensor product and $t>0$. There exists $C\ge0$
    such that, for every $k\ge0$, every label set $J$, not necessarily finite, every
    family of partitions $(\pi_j)_{j\in J}$, and every nonempty finite
    rooted weighted binary tree $\mathcal T$ with terminal labels in $J$,
    its root satisfies
    \begin{align}
        (k+2)^{-C}I
        \le\mathcal T\bigl(\widetilde A_k(\pi_j):j\bigr).
    \end{align}
    The exponent is chosen before the copy number, label set, partitions,
    tree and interpolation parameters. Repeated terminal labels, zero
    terminal weights, and interpolation parameters at either endpoint
    are allowed.
\end{theorem}

\begin{proof}\leanok
    \uses{thm:replica_symmetric_band_floor, thm:mean_tree_properties}
    Choose $C$ from Theorem~\ref{thm:replica_symmetric_band_floor} and put
    $c=(k+2)^{-C}>0$. Every terminal input is bounded below by the
    positive definite matrix $cI$. Monotonicity of the mean tree and its
    value on a constant input give
    \begin{align}
        cI=\mathcal T(cI:j)
        \le\mathcal T\bigl(\widetilde A_k(\pi_j):j\bigr).
    \end{align}
\end{proof}

% Source: 06-transport.tex, transport:means, lines 40--62;
% 07-comparators.tex, lower comparator argument, lines 603--621.
These are direct consequences of the replica floor and monotonicity of
finite mean trees used in the lower comparator argument
of \cite{OpenAI2026AreaLaw}.
